\hypertarget{ux6a21ux578bux538bux7f29}{%
\chapter{模型压缩}\label{ux6a21ux578bux538bux7f29}}

\hypertarget{ux6a21ux578bux91cfux5316}{%
\section{模型量化}\label{ux6a21ux578bux91cfux5316}}

\hypertarget{ux4e00ux6a21ux578bux91cfux5316ux7684ux6838ux5fc3ux601dux60f3}{%
\subsection{一、模型量化的核心思想}\label{ux4e00ux6a21ux578bux91cfux5316ux7684ux6838ux5fc3ux601dux60f3}}

由于大模型的计算、存储开销巨大，计算需要的时间较长，在大模型推理部署时，我们考虑降低模型中数值的精度，以使模型无需在专业的GPU上就可以运行。一般我们会将模型的32位浮点数转换为8位整数（有时还会选择转换为4位整数）格式。

值得注意的是，由于模型训练过程需要精确的参数进行梯度更新，模型训练时不应量化。对于物理模拟等对数值精度要求高的领域，也不宜运用模型量化。

\hypertarget{ux4e8c8-bitux91cfux5316}{%
\subsection{二、8-bit量化}\label{ux4e8c8-bitux91cfux5316}}

8-bit量化即量化为{[}-128,127{]}范围内的整数，共有256个可能的取值。

\hypertarget{ux5bf9ux79f0ux91cfux5316}{%
\subsubsection{1. 对称量化}\label{ux5bf9ux79f0ux91cfux5316}}

在这个示例张量中，绝对值最大的数是 \(10.8\)。为了保持对称量化，浮点有效范围扩展为 \([-10.8, 10.8]\)，再映射到 INT8 的对称区间。

\begin{longtable}[]{@{}rr@{}}
\toprule
浮点值 \(x\) & 对应 INT8 值 \(q\) \\
\midrule
\endhead
\(-7.59\) & \(-89\) \\
\(-4.57\) & \(-54\) \\
\(-1.95\) & \(-23\) \\
\(0\) & \(0\) \\
\(3.08\) & \(36\) \\
\(5.47\) & \(64\) \\
\(10.8\) & \(127\) \\
\bottomrule
\end{longtable}

寻找绝对值最大值：为了保持对称，必须以数据中绝对值最大的数为边界，定义新的量化范围（Clip Range）。

\[
\alpha = \max(|-7.59|, |10.8|) = 10.8
\]

因此，量化的有效浮点范围被扩展为 \([-10.8, 10.8]\)。

计算缩放因子 \(S\)：将上述浮点范围映射到 INT8 的对称区间 \([-127, 127]\)。

\[
S = \frac{10.8}{127} \approx 0.085039
\]

应用量化公式：由于 \(Z = 0\)，公式非常简洁：

\[
q = \mathrm{round}\left(\frac{x}{S}\right)
\]

\hypertarget{ux975eux5bf9ux79f0ux91cfux5316}{%
\subsubsection{2. 非对称量化}\label{ux975eux5bf9ux79f0ux91cfux5316}}

非对称量化直接使用原始浮点范围 \([-7.59, 10.8]\)，其中最大值为 \(10.8\)，最小值为 \(-7.59\)。

\begin{longtable}[]{@{}rr@{}}
\toprule
浮点值 \(x\) & 对应 INT8 值 \(q\) \\
\midrule
\endhead
\(-7.59\) & \(-128\) \\
\(-4.57\) & \(-86\) \\
\(-1.95\) & \(-50\) \\
\(0\) & \(-23\) \\
\(3.08\) & \(20\) \\
\(5.47\) & \(53\) \\
\(10.8\) & \(127\) \\
\bottomrule
\end{longtable}

确定范围：

\begin{itemize}
\tightlist
\item
  浮点范围：\(x_{\min} = -7.59, x_{\max} = 10.8\)
\item
  整数范围（INT8）：\(q_{\min} = -128, q_{\max} = 127\)
\end{itemize}

计算缩放因子 \(S\)：

\[
S = \frac{x_{\max} - x_{\min}}{q_{\max} - q_{\min}}
  = \frac{10.8 - (-7.59)}{127 - (-128)}
  = \frac{18.39}{255}
  \approx 0.072117
\]

注：非对称量化的 \(S\)（\(0.072\)）小于对称量化的 \(S\)（\(0.085\)），说明其量化粒度更细，精度更高。

计算零点 \(Z\)：计算浮点数 \(0.0\) 对应的整数值 \(Z\)。

\[
Z = \mathrm{round}\left(q_{\max} - \frac{x_{\max}}{S}\right)
\]

代入数值：

\[
Z = \mathrm{round}\left(127 - \frac{10.8}{0.072117}\right)
  \approx \mathrm{round}(127 - 149.75)
  \approx -23
\]

这意味着浮点数 \(0.0\) 映射为整数 \(-23\)。

应用量化公式：

\[
q = \mathrm{round}\left(\frac{x}{S} + Z\right)
\]

计算示例：

\begin{itemize}
\tightlist
\item
  对于 \(10.8\)：
\end{itemize}

\[
q = \mathrm{round}\left(\frac{10.8}{0.072117} - 23\right)
  = \mathrm{round}(149.75 - 23)
  = 127
\]

\begin{itemize}
\tightlist
\item
  对于 \(-7.59\)：
\end{itemize}

\[
q = \mathrm{round}\left(\frac{-7.59}{0.072117} - 23\right)
  = \mathrm{round}(-105.24 - 23)
  = -128
\]

优点：如示意图所示，原始范围 \([-7.59, 10.8]\) 完整覆盖了 \([-128, 127]\)，没有浪费任何量化层级，通常能保留更高的精度。

缺点：计算稍复杂，反量化或矩阵运算时需要额外处理 \(Z\) 的减法运算。

\hypertarget{ux4e094-bitux91cfux5316}{%
\subsection{三、4-bit量化}\label{ux4e094-bitux91cfux5316}}

4-bit量化将每个权重表示为0000-1111，共有16个值。

核心思想：我们不应该让量化点的间隔相等，而应该让每个量化点所代表的``概率分位数''的概率间隔相等。我们将(0,1)正态分布的累积分布函数等分为16个概率区间，取各区间中点对应的值，归一化映射到{[}-1,1{]}范围，作为NF4量化值。这些值和量化索引一一对应，如-1.0000对应0000，-0.6962对应0001，\ldots\ldots，1.0000对应1111.

量化步骤：

\begin{enumerate}
\def\labelenumi{\arabic{enumi}.}
\item
  计算缩放因子 \(S\)，将权重张量的范围映射到 NF4 量化表的 \([-1, 1]\) 范围。

  \begin{itemize}
  \tightlist
  \item
    计算张量中每个元素的绝对值：\([0.45, 0.88, 1.54, 0.12, 0.05, 0.99, 0.33, 2.10]\)
  \item
    找到最大值：\(2.10\)
  \item
    因此，缩放因子 \(S = 2.10\)
  \end{itemize}
\item
  归一化权重，用每个原始权重除以缩放因子 \(S\)，得到归一化后的权重 \(R_{norm} = \frac{R}{S}\)。

  \begin{itemize}
  \tightlist
  \item
    \(0.45 / 2.10 = 0.214\)
  \item
    \(-0.88 / 2.10 = -0.419\)
  \item
    \ldots\ldots{}
  \item
    由此得到归一化后的张量 \(R_{norm}\)：\([0.214, -0.419, 0.733, -0.057, 0.024, 0.471, -0.157, -1.0]\)
  \end{itemize}
\item
  映射到距离最近的NF4量化点，转换为对应的4位NF4索引。
\item
  存储结果，经过量化，最终需要存储两样东西：

  \begin{itemize}
  \tightlist
  \item
    4比特索引数组：\([10, 3, 14, 6, 8, 12, 5, 0]\)
  \item
    缩放因子 \(S\)：\(2.10\)（用 FP32 存储）
  \end{itemize}
\end{enumerate}

这个块的原始内存占用是 \(8 \times 32\) 位 \(= 256\) 位，现在是 \(8 \times 4\) 位 \(+ 32\) 位 \(= 64\) 位，内存压缩了 75\%。

反量化步骤（转换回FP32）：

\begin{enumerate}
\def\labelenumi{\arabic{enumi}.}
\item
  读取索引和缩放因子。

  \begin{itemize}
  \tightlist
  \item
    从内存中读取之前存储的索引数组和缩放因子 \(S = 2.10\)。
  \end{itemize}
\item
  从量化表恢复归一化值：使用索引法，在 NF4 量化表中查找对应的浮点值。

  \begin{itemize}
  \tightlist
  \item
    \(10 \to 0.1906\)
  \item
    \(3 \to -0.3949\)
  \item
    \ldots{}
  \item
    由此得到恢复的归一化张量 \(R'_{norm}\)：\([0.1906, -0.3949, 0.6717, -0.0911, 0.0062, 0.4045, -0.1848, -1.0]\)
  \end{itemize}
\item
  乘以缩放因子恢复权重：将恢复的归一化值乘以缩放因子 \(S\)，得到最终的反量化权重 \(R'\)。

  \begin{itemize}
  \tightlist
  \item
    \(R' = R'_{norm}S\)
  \item
    \(0.1906 \times 2.10 = 0.400\)
  \item
    \(-0.3949 \times 2.10 = -0.829\)
  \item
    \ldots{}
  \item
    最终得到 \(R'\)：\([0.400, -0.829, 1.411, -0.191, 0.013, 0.849, -0.388, -2.10]\)
  \end{itemize}
\end{enumerate}

\hypertarget{ux56dbux9759ux6001ux91cfux5316ux4e0eux52a8ux6001ux91cfux5316}{%
\subsection{四、静态量化与动态量化}\label{ux56dbux9759ux6001ux91cfux5316ux4e0eux52a8ux6001ux91cfux5316}}

\hypertarget{ux9759ux6001ux91cfux5316}{%
\subsubsection{1. 静态量化}\label{ux9759ux6001ux91cfux5316}}

（1）核心思想和步骤

静态量化是一种``全流程量化''策略。它的核心在于``离线（Offline）''完成所有繁重的工作。

校准：这是静态量化最关键的一步。因为还没有运行模型，我们不知道激活值（Activation，即每一层的输出）的范围。所以需要喂入少量真实的典型数据（校准集），让模型跑一遍，统计每一层激活值的min和max，从而提前计算出固定的S和Z。

转换：一旦有了S和Z，就将模型中所有的权重和激活值的处理逻辑全部写成INT8格式。

推理：运行时，模型不再需要做任何浮点运算转换，直接在整数域飞快奔跑。（两个整数相加可能超出INT8，此时会用INT32计算后转回INT8）

（2）优缺点和用途

优点：推理速度最快。

缺点：如果真实数据分布相对于校准集发生变化（NLP任务中尤其多见），S和Z可能不再适用；对于涉及Softmax等变换的场景（如Transformer），量化带来的微小差异可能在指数运算被急剧放大；不擅长处理``离群值''（如注意力机制中，在设置S和Z时，权重分布往往呈现大部分集中和一些离群值的状态，离群值往往起着类似``注意力锚点''的关键作用，但如果照顾离群值，则大部分值精度急剧下降）。

用途：CNN、边缘设备等。

\hypertarget{ux52a8ux6001ux91cfux5316}{%
\subsubsection{2. 动态量化}\label{ux52a8ux6001ux91cfux5316}}

（1）核心思想和步骤

动态量化是一种``折中''策略，核心是``权重静态，激活动态''，准备阶段只把模型的权重提前量化成 INT8。推理时：

层间传输：上一层的输出是FP32，下一层的输入也是FP32。

计算前：在进入矩阵乘法（MatMul）核心计算之前，激活值是FP32。系统需要读取这些 FP32 数据，计算统计量，然后将其转换为INT8。

计算时：核心的矩阵乘法是用 INT8 完成的（得到 INT32 结果）。

计算后：乘以 Scale 因子，将结果反量化回FP32。

非线性层：LayerNorm、GELU/SiLU等激活函数层通常直接在FP32上运行，不进行量化。

（2）优缺点和应用

优点：虽然不如静态量化快，但对于LLM这种内存受限的模型，加载权重的时间远多于计算时间。动态量化将权重体积缩小了4倍，意味着从内存加载权重的速度快了4倍，是性能提升的主要来源，与静态量化间的时间差相比而言很小。

缺点：需要实时分析当前激活值的范围，每次推理都要做一次min()、max()和div() 运算，并且要逐元素将FP32转为INT8，增加CPU/GPU的负担；数据在INT8核心和FP32显存之间频繁转换，增加了指令延迟。

应用：对于 Transformer/LLM（参数巨大，计算相对稀疏），动态量化通常能带来1.5x到2x的推理加速。因为``加载权重的收益''远大于``动态计算量化参数的开销''；对于 CNN（如 ResNet，计算密集型），动态量化效果往往不如静态量化，甚至可能比纯FP32慢，因为 CNN 的权重复用率高，计算量大，频繁的动态转换开销占比过高。

\hypertarget{ux53c2ux8003ux6587ux732e-86}{%
\subsection{参考文献}\label{ux53c2ux8003ux6587ux732e-86}}

\begin{itemize}
\tightlist
\item
  Jacob, B., Kligys, S., Chen, B., et al.~(2018). \href{https://arxiv.org/abs/1712.05877}{Quantization and Training of Neural Networks for Efficient Integer-Arithmetic-Only Inference}. CVPR 2018.
\item
  Dettmers, T., Lewis, M., Belkada, Y., \& Zettlemoyer, L. (2022). \href{https://arxiv.org/abs/2208.07339}{LLM.int8(): 8-bit Matrix Multiplication for Transformers at Scale}. NeurIPS 2022.
\item
  Dettmers, T., Pagnoni, A., Holtzman, A., \& Zettlemoyer, L. (2023). \href{https://arxiv.org/abs/2305.14314}{QLoRA: Efficient Finetuning of Quantized LLMs}. NeurIPS 2023.
\end{itemize}

\clearpage

\hypertarget{ux6a21ux578bux84b8ux998f}{%
\section{模型蒸馏}\label{ux6a21ux578bux84b8ux998f}}

\hypertarget{ux4e00ux6838ux5fc3ux601dux60f3-7}{%
\subsection{一、核心思想}\label{ux4e00ux6838ux5fc3ux601dux60f3-7}}

在深度学习中，我们通常训练一个庞大、复杂的模型来获得高准确率。但在实际部署（如手机端、IoT 设备）时，由于算力和内存限制，无法直接使用这个大模型。我们让一个小模型（Student）通过模仿一个大模型（Teacher）的行为，来学习大模型中蕴含的``暗知识''，从而在保持较小参数规模的同时，尽可能达到接近大模型的性能。下面引入两个概念：

Hard Target（硬标签）：老师告诉学生：答案是``A''。这对应传统的交叉熵损失函数，即训练时专注增加输出 A 的概率，不关心输出 B、C 的概率。

Soft Target（软标签）：老师告诉学生：``答案选 A，但 B 也有点道理，C 是大错特错的。''这种包含概率分布的信息（如 A=0.7, B=0.2, C=0.1）蕴含了类别之间的相似度关系，即大模型内的``暗知识''。

模型蒸馏会同时学习这两种标签。（不可只学习软标签，那样会导致输出正确答案的概率下降）

\hypertarget{ux4e8cux6570ux5b66ux8868ux8fbe-2}{%
\subsection{二、数学表达}\label{ux4e8cux6570ux5b66ux8868ux8fbe-2}}

\hypertarget{ux5e26ux6e29ux5ea6ux7684-softmax}{%
\subsubsection{1. 带温度的 Softmax}\label{ux5e26ux6e29ux5ea6ux7684-softmax}}

在标准的分类任务中，神经网络输出 logits \(z_i\)，经过 Softmax 得到概率 \(q_i\)。在蒸馏中，我们引入超参数 \(T\)：

\[
q_i = \frac{\exp(z_i / T)}{\sum_j \exp(z_j / T)}
\]

\begin{itemize}
\tightlist
\item
  当 \(T = 1\) 时，这就是标准的 Softmax。
\item
  当 \(T > 1\) 时，概率分布变得更``平滑''（Softer）。这使得原本接近 0 的错误类别的概率被放大，从而暴露出模型学到的类别间的隐含关系（例如：在识别``狗''时，模型认为像``猫''的概率比像``汽车''的概率大）。
\end{itemize}

\hypertarget{ux635fux5931ux51fdux6570loss-function}{%
\subsubsection{2. 损失函数（Loss Function）}\label{ux635fux5931ux51fdux6570loss-function}}

学生模型的总损失函数 \(L\) 通常由两部分组成：

\[
L = \alpha L_{KD} + (1 - \alpha)L_{CE}
\]

其中：

\begin{itemize}
\tightlist
\item
  \(L_{KD}\)（蒸馏损失/软损失）：衡量学生模型产生的软概率分布与教师模型产生的软概率分布之间的差异。通常使用 KL 散度（Kullback-Leibler Divergence）。注意，计算此时使用的是高温 \(T\)。
\item
  \(L_{CE}\)（学生损失/硬损失）：衡量学生模型的输出与真实标签（Ground Truth）之间的差异。通常使用标准的交叉熵损失（Cross-Entropy Loss）。注意，计算此时使用的是 \(T = 1\)。
\item
  \(\alpha\)：平衡两个损失权重的超参数。
\end{itemize}

注意这里的 KL 散度是 forward KL，因为：

KL 散度的标准数学定义为：

\[
D_{KL}(P \Vert Q) = \mathbb{E}_{x \sim P}\left[\log \frac{P(x)}{Q(x)}\right]
\]

请注意公式中的期望符号 \(\mathbb{E}_{x \sim P}\)。这意味着我们要从分布 \(P\) 中采样数据来计算这个积分或求和。

在机器学习的语境下，谁负责在物理世界（或计算图中）生成采样数据，谁就必须排在 KL 散度的前面。

\hypertarget{ux6570ux5b66ux63a8ux5bfcux4e3aux4ec0ux4e48ux9ad8ux6e29ux6709ux6548}{%
\subsubsection{3. 数学推导：为什么高温有效？}\label{ux6570ux5b66ux63a8ux5bfcux4e3aux4ec0ux4e48ux9ad8ux6e29ux6709ux6548}}

假设 logits \(z_i\) 相比于 \(T\) 很小（即 \(z_i / T \ll 1\)），我们可以对 Softmax 函数进行泰勒展开。Softmax 输出 \(q_i\) 大约为：

\[
q_i \approx \frac{1 + z_i / T}{N + \sum_j z_j / T}
\]

假设 logits 均值为 0（即 \(\sum_j z_j = 0\)），则简化为：

\[
q_i \approx \frac{1 + z_i / T}{N}
\]

假设我们使用交叉熵损失函数（或 KL 散度）进行蒸馏，其针对学生模型 Logit \(z_i^{(Student)}\) 的梯度通常形式为：

\[
\frac{\partial L}{\partial z_i^{(Student)}} = \frac{1}{T}\left(q_i^{(Student)} - q_i^{(Teacher)}\right)
\]

注：这里的 \(1/T\) 来自于链式法则中对 \(z_i / T\) 求导产生的系数。

现在，我们将上面的高温概率近似代入 \(q_i\)：

\[
\begin{aligned}
\frac{\partial L}{\partial z_i^{(Student)}} &\approx \frac{1}{T}\left[\left(\frac{1 + z_i^{(Student)} / T}{N}\right) - \left(\frac{1 + z_i^{(Teacher)} / T}{N}\right)\right] \\
&= \frac{1}{NT^2}\left(z_i^{(Student)} - z_i^{(Teacher)}\right)
\end{aligned}
\]

忽略常数系数 \(\frac{1}{NT^2}\)，我们发现梯度正比于：

\[
\frac{\partial L}{\partial z_i} \propto z_i^{(Student)} - z_i^{(Teacher)}
\]

回顾均方误差（MSE）的梯度：如果我们直接定义损失函数为 Logits 之间的 MSE：

\[
L_{MSE} = \frac{1}{2}\left(z^{(Student)} - z^{(Teacher)}\right)^2
\]

其梯度正是：

\[
\frac{\partial L_{MSE}}{\partial z^{(Student)}} = z^{(Student)} - z^{(Teacher)}
\]

这就证明了：在高温极限下，最小化 KL 散度（或软交叉熵）等价于最小化 Logits 之间的均方误差。

当然，实际上我们一般取 \(T\) 为 2\textasciitilde20，平衡 Hard Target 的``关注最大值''和 MSE 的``平衡关注各方向分布''。从实际意义上来说：

在正常温度下，Softmax 会把最大的 Logit 变成接近 1 的概率，而把其他 Logits 压缩成接近 0。

Teacher Logits: \texttt{{[}10,\ 5,\ 2{]}} -\textgreater{} Softmax: \texttt{{[}0.99,\ 0.009,\ 0.001{]}}

Teacher Logits: \texttt{{[}10,\ 8,\ 8{]}} -\textgreater{} Softmax: \texttt{{[}0.8,\ 0.1,\ 0.1{]}}

如果只看概率，很多细微的差异（比如 5 和 2 的区别）被``压扁''了，学生模型很难学到。

而 Logits 直接反映了模型对各个类别的原始置信度。

Logits \texttt{{[}10,\ 5,\ 2{]}} 明确告诉学生：虽然类别 1 最强，但类别 2 比类别 3 强得多（强 3 个单位），这种信息，就是暗知识（Dark Knowledge）。

上面模型蒸馏的数学表达，事实上正是``匹配 Logits''思想的体现。它的结果是：

\begin{itemize}
\tightlist
\item
  方向对齐：蒸馏过程要求学生生成的 Logits 向量 \(z_s\) 和教师的 Logits 向量 \(z_t\) 指向同一个方向。这意味着学生不仅要学会``谁是第一名''，还要学会``第二名、第三名\ldots\ldots 第 N 名''的相对排序。
\item
  相对大小对齐：\(z_i - z_j\) 代表了类别 \(i\) 和类别 \(j\) 的相对差距（Log-odds）。Logits Matching 强迫学生模型复刻这种相对差距。
\end{itemize}

例如，如果老师认为``猫''比``汽车''更像``狗''10 倍，学生也必须认为``猫''比``汽车''更像``狗''10 倍。

\hypertarget{ux4e09ux6a21ux578bux84b8ux998fux7684ux591aux79cdux65b9ux6cd5}{%
\subsection{三、模型蒸馏的多种方法}\label{ux4e09ux6a21ux578bux84b8ux998fux7684ux591aux79cdux65b9ux6cd5}}

\begin{enumerate}
\def\labelenumi{\arabic{enumi}.}
\item
  基于响应：利用教师模型最后一层的输出，这是最经典的方法（如上文所述）。
\item
  基于特征：不仅学习输出，还强制学生模型的中间层特征图拟合教师模型的中间层。通常需要引入额外的卷积层来匹配特征图的维度。
\end{enumerate}

损失函数：

\[
L_{\mathrm{Feature}} = \left\Vert F_{\mathrm{Teacher}}(x) - W_{\mathrm{reg}}\left(F_{\mathrm{Student}}(x)\right) \right\Vert^2
\]

\begin{enumerate}
\def\labelenumi{\arabic{enumi}.}
\setcounter{enumi}{2}
\tightlist
\item
  基于样本关系：
\end{enumerate}

对于同一批 \(N\) 个样本，分别获取它们在教师模型的特征空间中的特征向量（一般为其中一层的向量），以及在学生模型的特征空间中的特征向量。由于教师模型和学生模型特征向量维度不同，无法直接计算损失，故我们考虑计算向量两两之间的``相似度''（标量）。

我们对教师模型的这些特征向量两两计算余弦相似度或欧氏距离，得到 \(N \times N\) 的``相似度矩阵''，对学生模型也得到这样的矩阵。如果用的是欧氏距离，为了消除维度差异，我们还会加归一化因子。我们的目标是让这两个矩阵长得像，即结构一致性：

\[
L_{\mathrm{Relation}} = \sum_{i=1}^{N} \sum_{j=1}^{N} L_\delta\left(\mathcal{R}_{ij}^{(T)}, \mathcal{R}_{ij}^{(S)}\right)
\]

其中 \(L_\delta\) 可以是 Huber Loss 或 \(L_2\) Loss。

Huber Loss 是一种鲁棒回归损失函数，它结合了均方误差（MSE/L2）和平均绝对误差（MAE/L1）的优点。

简单来说，它的设计哲学是：``小错误用平方惩罚（为了平滑收敛），大错误用线性惩罚（为了防止梯度爆炸）。''

Huber Loss \(L_\delta(y, f(x))\) 对于预测误差 \(a = y - f(x)\) 的定义如下：

\[
L_\delta(a) =
\begin{cases}
\frac{1}{2}a^2, & \text{for } |a| \le \delta \\
\delta \cdot \left(|a| - \frac{1}{2}\delta\right), & \text{for } |a| > \delta
\end{cases}
\]

其中：

\begin{itemize}
\tightlist
\item
  \(a\) 是残差（真实值与预测值的差）。
\item
  \(\delta\)（delta）是一个超参数，作为``分界点''。
\end{itemize}

当误差很小（\(|a| \le \delta\)）时：它是一个抛物线（像 MSE）。

\begin{itemize}
\tightlist
\item
  优点：在 0 点附近是可导的，梯度会随着误差变小而变小，有利于模型在最后阶段精细收敛，不会像 MAE 那样在 0 点附近震荡。
\end{itemize}

当误差很大（\(|a| > \delta\)）时：它变成了直线（像 MAE）。

\begin{itemize}
\tightlist
\item
  优点：它的梯度是常数 \(\delta\)（或者 \(-\delta\)）。这意味着即使遇到极端的离群点（Outliers），产生的梯度也不会无穷大，避免了 MSE 那种``因为一个离群点导致整个模型参数被带偏''的情况。
\end{itemize}

对噪声鲁棒：教师模型并不是完美的，它产生的关系矩阵可能包含噪声。如果某两个样本，教师认为距离是 100（可能是异常值），而学生认为是 1，由于 MSE 是平方增长，Loss 会变得巨大（\(99^2 \approx 9800\)），这会导致学生模型的梯度剧烈波动。使用 Huber Loss，这部分的惩罚是线性的，相对温和。

\hypertarget{ux56dbon-policy-ux84b8ux998f}{%
\subsection{四、On-policy 蒸馏}\label{ux56dbon-policy-ux84b8ux998f}}

传统的模型蒸馏是 Off-policy 的，教师模型自回归生成，学生模型则每次基于问题和教师模型的第 \(1 \sim i-1\) 个 token 生成第 \(i\) 个 token，让它的输出分布接近教师模型的输出分布。这就类似于描红，学生拿着半透明的描红纸盖在上面。教师写了前三笔，学生就顺着前三笔的笔势去``猜''第四笔该落在哪。学生全程都被教师的笔迹限制在一条正确的路径上。

既然学生模型在训练时从来没有自己自回归生成，那么当它被部署到真实环境（Test Time）中时，一旦没有了大师的字帖打底，学生必须自己从头写到尾。如果它在第 3 步不小心写歪了一点点，到了第 4 步时，它面对的上下文就是一个包含错误的、自己从未在训练中见过的陌生前缀。由于它在训练时只见过完美的路径，完全不知道怎么在错误路径上纠偏，接下来的生成就会彻底崩溃。

On-policy 蒸馏则较好地解决了这个问题：让学生模型自己在环境中自由探索，自回归生成完整的轨迹 \(y\)。即便这条轨迹包含错误、绕路或表述不佳，也都保留下来。把这条``充满学生个人风格（甚至错误）''的轨迹喂给教师大模型。教师大模型就像批改作业一样，针对这条轨迹的每一个时间步，给出全词表上的标准概率分布（Logits）。学生模型再根据这个分布来更新自己的参数。学生模型完全在自己的能力边界和状态空间内进行探索。它学到的是``如果在我的水平下走到了这一步，大模型会建议我下一步怎么救场或继续''。由于通过学生的轨迹采样，因此这里使用 Reverse KL（反向 KL 散度）学习：

\[
D_{KL}(\pi_{\mathrm{student}} \Vert \pi_{\mathrm{teacher}}) = \mathbb{E}_{y \sim \pi_{\mathrm{student}}}\left[\log \frac{\pi_{\mathrm{student}}(y|x)}{\pi_{\mathrm{teacher}}(y|x)}\right]
\]

这样可以使模型分布稳定在教师``圈定''的合理范围内，保证输出质量。

\hypertarget{ux53c2ux8003ux6587ux732e-87}{%
\subsection{参考文献}\label{ux53c2ux8003ux6587ux732e-87}}

\begin{itemize}
\tightlist
\item
  Hinton, G., Vinyals, O., \& Dean, J. (2015). \href{https://arxiv.org/abs/1503.02531}{Distilling the Knowledge in a Neural Network}. arXiv:1503.02531.
\item
  Agarwal, R., Vieillard, N., Zhou, Y., et al.~(2024). \href{https://arxiv.org/abs/2306.13649}{On-Policy Distillation of Language Models: Learning from Self-Generated Mistakes}. arXiv:2306.13649.
\item
  Song, M., \& Zheng, M. (2026). \href{https://arxiv.org/abs/2604.00626}{A Survey of On-Policy Distillation for Large Language Models}. arXiv:2604.00626.
\end{itemize}

\clearpage

\hypertarget{ux6a21ux578bux526aux679d}{%
\section{模型剪枝}\label{ux6a21ux578bux526aux679d}}

模型剪枝，即移除模型中``不重要''的参数（权重）或结构。分为两类：

\hypertarget{ux4e00ux975eux7ed3ux6784ux5316ux526aux679d}{%
\subsection{一、非结构化剪枝}\label{ux4e00ux975eux7ed3ux6784ux5316ux526aux679d}}

移除单个权重（将其设为0）。由于可以根据模型权重动态调压缩率高，但通常需要专用硬件/库支持稀疏计算才能获得加速。

\hypertarget{ux4e8cux7ed3ux6784ux5316ux526aux679d}{%
\subsection{二、结构化剪枝}\label{ux4e8cux7ed3ux6784ux5316ux526aux679d}}

移除整个结构单元，如滤波器（通道）、神经元、层。更易于在通用硬件上获得加速，但压缩率可能稍低。

流程通常包括：训练一个大模型 -\textgreater{} 评估参数重要性（如根据绝对值大小或梯度） -\textgreater{} 移除不重要的参数 -\textgreater{} 对剪枝后的模型进行微调。现代方法常采用迭代剪枝。

\hypertarget{ux53c2ux8003ux6587ux732e-88}{%
\subsection{参考文献}\label{ux53c2ux8003ux6587ux732e-88}}

\begin{itemize}
\tightlist
\item
  Han, S., Pool, J., Tran, J., \& Dally, W. (2015). \href{https://papers.nips.cc/paper_files/paper/2015/hash/ae0eb3eed39d2bcef4622b2499a05fe6-Abstract.html}{Learning both Weights and Connections for Efficient Neural Network}. NeurIPS 2015.
\item
  Li, H., Kadav, A., Durdanovic, I., Samet, H., \& Graf, H. P. (2017). \href{https://openreview.net/forum?id=rJqFGTslg}{Pruning Filters for Efficient ConvNets}. ICLR 2017.
\item
  Frankle, J., \& Carbin, M. (2019). \href{https://openreview.net/forum?id=rJl-b3RcF7}{The Lottery Ticket Hypothesis: Finding Sparse, Trainable Neural Networks}. ICLR 2019.
\end{itemize}

\clearpage

\hypertarget{ux9ad8ux7ef4ux5411ux91cfux7684ux91cfux5316ux8bbaux6587}{%
\section{高维向量的量化（论文）}\label{ux9ad8ux7ef4ux5411ux91cfux7684ux91cfux5316ux8bbaux6587}}

\begin{quote}
本文是论文阅读笔记，内容代表对应论文方法或作者理解，不应直接视为领域共识或工程最佳实践。
\end{quote}

\hypertarget{ux4e00ux95eeux9898ux80ccux666f}{%
\subsection{一、问题背景}\label{ux4e00ux95eeux9898ux80ccux666f}}

在人工智能和深度学习领域，高效处理高维向量至关重要。例如：

\begin{itemize}
\tightlist
\item
  \textbf{大语言模型（LLM）的 KV Cache}：自回归模型在生成长文本时，需要存储大量历史 Token 的键值对（Key-Value Cache）。随着上下文变长，内存消耗和访存延迟会显著增加。通过量化压缩这些向量，可以大幅降低推理成本。
\item
  \textbf{向量数据库与近似最近邻搜索（ANNS）}：在检索增强生成（RAG）和信息检索中，需要快速计算查询向量与数据库向量的内积或余弦相似度。
\end{itemize}

\textbf{现有方法的局限性}：目前的量化算法往往面临两难选择。要么计算速度慢、缺乏硬件加速器（如 GPU）的兼容性，无法满足在线（Online）实时量化的需求；要么在给定位宽（Bit-width）下无法达到最优的失真下界。此外，许多现有方法依赖于离线数据（Data-dependent），需要耗时的预处理和微调。

下面给出了一些学界就高维向量量化算法的研究。

\hypertarget{ux4e8cux9488ux5bf9ux5747ux65b9ux8befux5deeux4f18ux5316ux7684ux7b97ux6cd5}{%
\subsection{二、针对均方误差优化的算法}\label{ux4e8cux9488ux5bf9ux5747ux65b9ux8befux5deeux4f18ux5316ux7684ux7b97ux6cd5}}

\hypertarget{ux9488ux5bf9ux5747ux65b9ux8befux5deeux4f18ux5316ux7684-mathrmturboquant_mathrmmse}{%
\subsubsection{\texorpdfstring{1. 针对均方误差优化的 \(\mathrm{TurboQuant}_{\mathrm{mse}}\)}{1. 针对均方误差优化的 \textbackslash mathrm\{TurboQuant\}\_\{\textbackslash mathrm\{mse\}\}}}\label{ux9488ux5bf9ux5747ux65b9ux8befux5deeux4f18ux5316ux7684-mathrmturboquant_mathrmmse}}

该算法旨在最小化量化前后的均方误差（MSE）。其核心理论依据是：将输入向量进行随机旋转后，其各个坐标在数学上会服从 Beta 分布（在高维下渐近于高斯分布），且不同坐标之间几乎相互独立。因此，可以对每个坐标独立应用最优的标量量化（Lloyd-Max 算法）。

\(\mathrm{TurboQuant}_{\mathrm{mse}}\) 的工作流：

\begin{enumerate}
\def\labelenumi{\arabic{enumi}.}
\tightlist
\item
  \textbf{初始化（预计算）}：生成一个全局的随机旋转矩阵 \(\Pi \in \mathbb{R}^{d \times d}\)。针对 Beta 分布求解一维连续 K-Means 问题，预先计算并存储最优的聚类中心（Centroids）代码本 \(c_1,c_2,\ldots,c_{2^b}\)。
\item
  \textbf{量化过程（QUANT）}：

  \begin{itemize}
  \tightlist
  \item
    接收输入向量 \(x \in \mathbb{R}^d\)。
  \item
    应用随机旋转矩阵将其转换为 \(y = \Pi \cdot x\)。
  \item
    对于旋转后向量的每一个坐标 \(y_j\)，在代码本中找到距离最近的聚类中心 \(c_k\)，并记录其索引。
  \item
    输出这些由 \(b\) 个比特表示的索引向量。
  \end{itemize}
\item
  \textbf{反量化过程（DEQUANT）}：

  \begin{itemize}
  \tightlist
  \item
    接收索引向量，从代码本中查找对应的聚类中心，重建向量 \(\tilde{y}\)。
  \item
    应用逆旋转矩阵进行恢复：
  \end{itemize}
\end{enumerate}

\[
\tilde{x} = \Pi^\top \cdot \tilde{y}.
\]

\hypertarget{ux4e09ux9488ux5bf9ux5185ux79efux65e0ux504fux4f30ux8ba1ux7684ux7b97ux6cd5}{%
\subsection{三、针对内积无偏估计的算法}\label{ux4e09ux9488ux5bf9ux5185ux79efux65e0ux504fux4f30ux8ba1ux7684ux7b97ux6cd5}}

MSE 最优的量化器在估计内积时会引入偏差，这在低位宽时尤其明显。故需要针对内积无偏估计的算法。

\begin{enumerate}
\def\labelenumi{\arabic{enumi}.}
\tightlist
\item
  \textbf{初始化}：实例化一个位宽为 \(b-1\) 的 \(\mathrm{QUANT}_{\mathrm{mse}}\) 量化器，并生成一个元素服从标准正态分布的随机投影矩阵 \(S \in \mathbb{R}^{d \times d}\)。
\item
  \textbf{量化过程（QUANT）}：

  \begin{itemize}
  \tightlist
  \item
    使用 \(b-1\) 位的 \(\mathrm{QUANT}_{\mathrm{mse}}\) 对输入向量 \(x\) 进行量化，得到索引 \(\mathrm{idx}\)。
  \item
    通过反量化计算残差向量：
  \end{itemize}
\end{enumerate}

\[
r = x - \mathrm{DEQUANT}_{\mathrm{mse}}(\mathrm{idx}).
\]

\begin{itemize}
\tightlist
\item
  对残差向量应用 1-bit QJL 变换：
\end{itemize}

\[
\mathrm{qjl} = \mathrm{sign}(S \cdot r).
\]

\begin{itemize}
\tightlist
\item
  输出元组 \((\mathrm{idx}, \mathrm{qjl}, \lVert r \rVert_2)\)，包含 MSE 索引、残差的 1-bit 符号向量和残差的 L2 范数。
\end{itemize}

\begin{enumerate}
\def\labelenumi{\arabic{enumi}.}
\setcounter{enumi}{2}
\tightlist
\item
  \textbf{反量化过程（DEQUANT）}：

  \begin{itemize}
  \tightlist
  \item
    首先通过 \(\mathrm{DEQUANT}_{\mathrm{mse}}(\mathrm{idx})\) 恢复出基础向量 \(\tilde{x}_{\mathrm{mse}}\)。
  \item
    通过 QJL 的逆映射恢复残差向量：
  \end{itemize}
\end{enumerate}

\[
\tilde{x}_{\mathrm{qjl}} = \frac{\sqrt{\pi/2}}{d} \cdot \lVert r \rVert_2 \cdot S^\top \cdot \mathrm{qjl}.
\]

\begin{itemize}
\tightlist
\item
  将两者相加得到最终重建向量：
\end{itemize}

\[
\tilde{x} = \tilde{x}_{\mathrm{mse}} + \tilde{x}_{\mathrm{qjl}}.
\]

\hypertarget{ux53c2ux8003ux6587ux732e-89}{%
\subsection{参考文献}\label{ux53c2ux8003ux6587ux732e-89}}

\begin{itemize}
\tightlist
\item
  Zandieh, A., Daliri, M., \& Han, I. (2024). \href{https://arxiv.org/abs/2406.03482}{QJL: 1-Bit Quantized JL Transform for KV Cache Quantization with Zero Overhead}. arXiv:2406.03482.
\item
  Vargaftik, S., Ben Basat, R., Portnoy, A., Mendelson, G., \& Mitzenmacher, M. (2021). DRIVE: One-bit Distributed Mean Estimation. NeurIPS 2021.
\item
  Vargaftik, S., et al.~(2022). EDEN: Communication-Efficient and Robust Distributed Mean Estimation for Federated Learning. ICML 2022.
\item
  Wu, X., et al.~(2024). RaBitQ: Quantizing High-Dimensional Vectors with a Theoretical Error Bound for Approximate Nearest Neighbor Search. SIGMOD 2024.
\item
  Zandieh, A., Daliri, M., Hadian, M., \& Mirrokni, V. (2025). \href{https://arxiv.org/abs/2504.19874}{TurboQuant: Online Vector Quantization with Near-optimal Distortion Rate}. arXiv:2504.19874.
\end{itemize}

\clearpage
